\documentclass[12pt]{article}

% фонтови и језик
% fontspec docs: shorturl.at/ouI26
\usepackage{fontspec}

% polyglossia docs: shorturl.at/gELX9
\usepackage{polyglossia}
\newfontfamily\cyrillicfont{Times New Roman} % може се заменити компатибилним фонтом
\newfontfamily\cyrillicfontsf{Arial}         % може се заменити компатибилним фонтом
\newfontfamily\cyrillicfonttt{Courier New}   % може се заменити компатибилним фонтом
\setmainlanguage[Script=cyrillic]{serbian}
\setotherlanguage{english}

% вербатим линкови (веб адресе, имејлови, релативне адресе, итд.)
% url docs: shorturl.at/iktM4
\usepackage{url}

% хиперлинкови
% hyperref docs: shorturl.at/fmHPW
\usepackage{hyperref}

% приказивање математичких израза
% amsmath docs: shorturl.at/avJU3
\usepackage{amsmath}

% напредни пакет за исцрватање графика коришћењем команде \includegraphics
% graphicx docs: shorturl.at/diHUW
\usepackage{graphicx}
\graphicspath{{images/}}    % коренски директоријум слика, сваки пут када се користи includegraphics команда путања која се задаје треба да буде релативна у односу на овај директоријум

% рад са прилозима
% appendix docs: shorturl.at/uwEL1
\usepackage{appendix}

% подешавања маргина
\usepackage{vmargin}
\setmarginsrb{3 cm}{2.5 cm}{3 cm}{2.5 cm}{1 cm}{1.5 cm}{1 cm}{1.5 cm}

% исцртавање програмског кода
\usepackage{minted}

% подешавање размака линија текста
\usepackage{setspace}

% цртање табела
\usepackage{booktabs}

% додатни макрои за израду рада
\usepackage{rsvp}


\title{Наслов рада}                         % ТОДО изменити
\author{Петар Трифуновић}                  % ТОДО изменити
\newcommand{\studentindex}{E2 4/2021}     % ТОДО изменити

\usepackage{fancyhdr}

\makeatletter
\let\thetitle\@title
\let\theauthor\@author
\let\thedate\@date
\let\theindex\studentindex
\makeatother

\pagestyle{fancy}
\fancyhf{}
\rhead{\theauthor}
\lhead{\thetitle}
\cfoot{\thepage}


\begin{document}
\begin{titlepage}
	\centering
    \vspace*{0.5 cm}
    \includegraphics[scale = 0.75]{ftn-logo.eps}\\[1.0 cm]	                % University Logo
    \textsc{\LARGE Факултет техничких наука}\\[0.5 cm]	                    % University Name
	\textsc{\Large Универзитет у Новом Саду}\\[1.0 cm]				        % Course Code
	\textsc{\large Рачунарски системи високих перформанси}\\[0.5 cm]		% Course Name
	\rule{\linewidth}{0.2 mm} \\[0.4 cm]
	{ \huge \bfseries \thetitle}\\
	\rule{\linewidth}{0.2 mm} \\[1.5 cm]
	
	\begin{minipage}{0.4\textwidth}
		\begin{flushleft} \large
			\emph{Аутор:}\\
			\theauthor
			\end{flushleft}
			\end{minipage}~
			\begin{minipage}{0.4\textwidth}
			\begin{flushright} \large
			\emph{Индекс:} \\
			\theindex								                        % Your Student Number
		\end{flushright}
	\end{minipage}\\[2.0 cm]
	
	{\large \thedate}\\[2 cm]
	\vfill
\end{titlepage}

\pagenumbering{roman}
\begin{abstract}
\doublespacing
Тема овог рада јесте дистрибуирано односно паралелно претраживање графа по ширини. Мотивација за писање рада јесте пре свега истраживање начина за убрзање релативно основних опереција над веома битном структуром података у рачунарству, односно графом. Такође, дубље истраживање проблема који се јавља у дистрибуираном и паралелном програмирању када се говори о претраживању графова. Наиме, због структуре која је, за недостатак бољих речи, непредвидива, тешко је добити прихватљиве перформансе при приступању нелинеарно складиштеним податацима. Кроз рад су анализиране методе за расподелу посла између процеса односно нити који врше саму претрагу, предности и мане оба приступа, као и постигнути резултати при њиховој примени. Закључено је да иако наизглед комплексно и мукотрпно, ипак је могуће постигнути неко убрзање кроз паралелни приступ у окружењима са дељеном меморијом. Нажалост, што се овог рада тиче, имплементација у оквиру окружења са дистрибуираном меморијом није показала позитивне резултате.
\end{abstract}
\pagebreak          % ТОДО попунити              
%%%%%%%%%%%%%%%%%%%%%%%%%%%%%%%%%%%%%%%%%%%%%%%%%%%%%%%%%%%%%%%%%%%%%%%%%%%%%%%%%%%%%%%%%

% садржај
\tableofcontents
\pagebreak

%%%%%%%%%%%%%%%%%%%%%%%%%%%%%%%%%%%%%%%%%%%%%%%%%%%%%%%%%%%%%%%%%%%%%%%%%%%%%%%%%%%%%%%%%

% листинг изворних кодова
\listoflistings
\pagebreak

%%%%%%%%%%%%%%%%%%%%%%%%%%%%%%%%%%%%%%%%%%%%%%%%%%%%%%%%%%%%%%%%%%%%%%%%%%%%%%%%%%%%%%%%%

% листинг фигура
\listoffigures
\pagebreak

%%%%%%%%%%%%%%%%%%%%%%%%%%%%%%%%%%%%%%%%%%%%%%%%%%%%%%%%%%%%%%%%%%%%%%%%%%%%%%%%%%%%%%%%%

% листинг табела
\listoftables
\pagebreak

%%%%%%%%%%%%%%%%%%%%%%%%%%%%%%%%%%%%%%%%%%%%%%%%%%%%%%%%%%%%%%%%%%%%%%%%%%%%%%%%%%%%%%%%%

\pagenumbering{arabic}
\section{Увод}
У уводу треба дати ширу мотивацију за проблем који решавате. Детаљније описати проблем и концептуално методологију решавања. Ту методологији ћете описати у детаље у засебном поглављу. Затим, ако већ постоје слична решења проблема који решавате, можете направити и кратак преглед постојећих решења уз навођење одговарајућих референци. Референце се задају у форми \textit{bibtex} уноса. У библиографији се појављују само ако их цитирате негде у тексту. Пример цитирања референце из библиографије: \cite{bibtex}.              % ТОДО попунити

\section{Назив првог поглавља тела рада}
Ово поглавље је намењено као демонстрација неких основних \LaTeX команди за рад са текстом. Поглавље треба искључити из финалне верзије вашег рада.

\subsection{Први поднаслов}
Уколико ћете имати само један поднаслов у оквиру поглавља, онда напишите све у оквиру поглавља без поднаслова.

\subsection{Други поднаслов}
Уколико имате \textbf{први поднаслов}, онда ћете сигурно имати \textit{најмање} још један. Можете додати произвољан број поднаслова додајући \texttt{subsection} команду.\\

Ово је други параграф у оквиру поднаслова, добијен је коришћењем \texttt{\textbackslash\textbackslash} команде на крају првог параграфа.

   % ТОДО направити нова поглавља по узору на poglavlje1 и poglavlje2
\section{Примери}
У овом поглављу су приказани примери исцртавања програмског кода, фигуре и табеле у оквиру документа. Сврха примера је да вам олакша рад са \LaTeX-ом уколико раније нисте радили са њим. Примери су само референтни и можете изабрати и друге пакете за исцртавање од оних који су коришћени у примерима.

Још примера цитирања референци \cite{kamran2101rv, yatsura2021metalearning}.


\subsection{Пример приказа програмског кода коришћењем \texttt{minted} пакета}
Коришћењем команди minted пакета, могуће је исцртати код као посебан блок или у склопу текста (\textit{inline} режим). Изворни код се може задати директно или учитати из датотеке.

\subsubsection{Исцртавање изворног кода из датотеке}
\inputpython{kodovi/helloworld.py}{Пример python кода}{code:helloworldpy}
Пример \ref{code:helloworldpy} Python изворног кода је добијен коришћењем \LaTeX команде \texttt{inputpython} дефинисане у датотеци \texttt{rsvp.sty}. Команда је омотач за \texttt{inputminted} команду са бојењем за Python програмски језик и нумерисање изворних кодова како би се могли референцирати и додати у листу свих изворних кодова. Имплементирани су још и омотачи за C и C++ програмске језике. По узору на имплементиране можете додати и своје омотаче. 

\begin{listing}
\inputminted{c}{kodovi/helloworld.c}
\caption{Пример \texttt{C} кода}
\label{code:helloworldc}
\end{listing}
Код за испис примера \ref{code:helloworldc} показује како можете исцртати изворни код у C програмском језику без коришћења омотача из \texttt{rsvp.sty} датотеке. Може бити корисно уколико желите финију контролу над исцртавањем кода коју омотачи не поджавају.

\subsubsection{Исцртавање блока изворног кода}
\begin{pythoncode}{Још један пример исцртавања \texttt{python} изворног кода}{code:hellowpyblock}
def hello_world():
    print("Hello world!")


if __name__ == '__main__':
    hello_world()
\end{pythoncode}

\begin{listing}
\begin{minted}{c}
#include <stdio.h>

void hello_world() {
    printf("Hello world!");
}

int main() {

    hello_world();

    return 0;
}
\end{minted}
\caption{Још један пример C изворног кода}
\label{code:hellowcblock}
\end{listing}

Изворни код \ref{code:hellowpyblock} је исцртан коришћењем \texttt{pythoncode} окружења, док је изворни код \ref{code:hellowcblock} исцртан без коришћења омотача.

\subsubsection{Исцртавање \textit{inline} изворног кода}
Код може бити исцртан и у оквиру текста са или без коришћења омотача. Нпр. \pythoninline{import os}, тј. \mintinline{c}{#include <stdio.h>}.

\subsection{Пример приказа фигуре коришћењем \texttt{graphicx} пакета}
Овде је приказано како исцртати и референцирати се на фигуру која садржи центрирану, нескалирану слику (нпр. слика \ref{fig:ftnlogo}).
\begin{figure}[H]
    \centering
    \includegraphics{images/ftn-logo.eps}
    \caption{ФТН лого}
    \label{fig:ftnlogo}
\end{figure}


\subsection{Пример приказа табеле коришћењем \texttt{tabular} пакета}
\begin{table}[H]
\centering
\begin{tabular}{@{}ccccc@{}}    % c значи да ће садржај ћелија табеле бити центриран
\toprule
колона 1 & колона 2 & колона 3 & колона 4 & колона 5 \\ \midrule\midrule
1 & 1 & 1 & 1 & 1 \\
1 & 1 & 1 & 1 & 1 \\
1 & 1 & 1 & 1 & 1 \\ \bottomrule
\end{tabular}
\caption{Пример табеле са заглављем}
\label{table:example1}
\end{table}

\section{Закључак}
Овде врло кратко подсетити на решаван проблем и написати закључке до којих сте дошли. Можете додати и будуће правце рада који нису имплементирани, а могли би допринети побољшању даље студије проблема.         % ТОДО попунити
\newpage

\begin{appendices}

\section{Први додатак}
Додаци су необавезни део рада. Овде се могу убацити сирови подаци на основу којих сте, на пример, нацртали графиконе приказане у телу рада, а којих је исувише да би се приказали у самом телу рада без одвлачења пажње читаоца са шире слике утицаја добијених резултата.

\section{Други додатак}
ТОДО опционо

\end{appendices}
            % ТОДО попунити уколико су вам потребни, додаци могу бити изостављени
\newpage

\bibliographystyle{plain}
\typeout{}
\bibliography{biblist}              % ТОДО попуњавати референцама

\end{document}